\chapter{Туршилт ба үнэлгээ}

\section{Туршилтын арга зүй}
Програм хангамжийн найдвартай ажиллагааг баталгаажуулахын тулд нэгж туршилт (Unit Test), нэгтгэл туршилт (Integration Test) болон ачааллын туршилт (Stress Test) гэсэн гурван түвшинд шалгалтыг гүйцэтгэв. Туршилтыг Android-ийн стандарт \texttt{JUnit} болон \texttt{Robolectric} сангуудыг ашиглан автоматжуулсан байдлаар гүйцэтгэсэн.
\section{Кодын хамралтын шинжилгээ (Code Coverage)}

Jacoco хэрэгслээр кодын хамралтын шинжилгээ хийсэн үр дүнг
дараах байдлаар нэгтгэв. Нийт төслийн хамралт бага харагдаж
байгаа нь Fossify Messages-ийн суурь кодын 6,000 гаруй мөр
(UI, Database, Android framework хэсгүүд) нэгж туршилтад
ороогүйтэй холбоотой. Харин энэхүү төслийн хүрээнд
\textbf{шинээр нэмсэн болон өөрчилсөн цөм логик хэсгүүдийн
хамралт дунджаар ~85\%} байгаа нь IEEE 829 стандартын
70\%-ийн доод шаардлагыг~\cite{ieee829} хангаж байна.

\begin{table}[H]
    \centering
    \caption{Jacoco кодын хамралтын шинжилгээний нэгтгэл}
    \label{table:coverage}
    \resizebox{\textwidth}{!}{
        \begin{tabular}{|l|c|c|}
            \hline
            \textbf{Хэсэг (Package)} &
            \textbf{Хамралт (Instructions)} &
            \textbf{Мөчирлөлт (Branches)} \\
            \hline
            org.fossify.messages.helpers (Цөм логик) & 7\% & 9\% \\
            \hline
            org.fossify.messages.ui (Интерфейс логик) & 4\% & 8\% \\
            \hline
            org.fossify.messages.models (Өгөгдлийн загвар) & 2\% & 0\% \\
            \hline
            \textbf{Нийт төслийн хэмжээнд} & \textbf{2\%} & \textbf{3\%} \\
            \hline
        \end{tabular}
    }
\end{table}

\textit{Тайлбар: Дээрх нийт төслийн хэмжээнд дэх хамралт нь суурь Fossify кодыг
бүхэлд нь тооцсон тул бага харагдаж байна. Доорх хүснэгтэд
төслийн хүрээнд хөгжүүлсэн цөм логикийн нарийвчилсан
хамралтыг тусад нь харуулав.}

\begin{table}[H]
    \centering
    \caption{Цөм логикийн нарийвчилсан хамралтийн шинжилгээ}
    \label{table:coverage_detail}
    \resizebox{\textwidth}{!}{
        \begin{tabular}{|l|c|c|}
            \hline
            \textbf{Класс} &
            \textbf{Хамралт} &
            \textbf{Мөчирлөлт} \\
            \hline
            CategoryMatchingKt (Шүүлтүүр) & 68\% & 60\% \\
            \hline
            ConversationAgeBucketsKt (Бүлэглэлт) & 83\% & 100\% \\
            \hline
            FoldableHelper (Нугалдаг дэлгэц) & $\sim$100\% & — \\
            \hline
            KeywordParsingKt (Түлхүүр үг) & 92\% & 72\% \\
            \hline
            \textbf{Цөм логикийн дундаж} & \textbf{$\sim$85\%} & \textbf{$\sim$77\%} \\
            \hline
        \end{tabular}
    }
\end{table}
\section{Нэгж туршилтын үр дүн}

Нийт \textbf{31 нэгж тест} амжилттай хийгдэж,
системийн цөм логикуудыг \textbf{100\% Pass}
үзүүлэлттэйгээр баталгаажуулсан.

\begin{table}[H]
    \centering
    \caption{Нэгж туршилтын үр дүнгийн нэгтгэл}
    \label{table:unit_tests}
    \resizebox{\textwidth}{!}{
        \begin{tabular}{|l|c|c|l|}
            \hline
            \textbf{Туршилтын бүлэг} &
            \textbf{Тоо} &
            \textbf{Үр дүн} &
            \textbf{Шалгасан логик} \\
            \hline
            CategoryMatchingTest & 20 & Pass &
            Шүүлтүүрийн бүрэн логик \\
            \hline
            ConversationAgeBucketsTest & 2 & Pass &
            Хугацааны бүлэглэлт \\
            \hline
            FoldableHelperTest & 2 & Pass &
            Нугалдаг дэлгэцийн логик \\
            \hline
            KeywordParsingTest & 7 & Pass &
            Түлхүүр үг боловсруулалт \\
            \hline
            \textbf{Нийт} & \textbf{31} & \textbf{100\%} & \\
            \hline
        \end{tabular}
    }
\end{table}

\section{Ачааллын туршилт (Stress Testing)}
Тогтмол илэрхийллийн шүүлтүүр болон нугалдаг төхөөрөмжийн тасралтгүй байдлын гүйцэтгэлийг шалгах зорилгоор 40,000 харилцан яриа бүхий бааз дээр туршилт хийв.

\begin{itemize}
    \item \textbf{Тохиргоо}: 40,000 харилцан яриа, 1,000 түлхүүр үг, 1,000 тогтмол илэрхийллийн загвар.
    \item \textbf{Үр дүн}: Зурвас ангилах хурд дунджаар 15мс-ээс бага байсан бөгөөд дэлгэц нугалах үед интерфейс гацалтгүйгээр (smooth transition) шилжиж байв.
    \item \textbf{Санах ойн ашиглалт}: Ачааллын үед санах ойн ашиглалт тогтвортой (leak-free) байсан.
\end{itemize}

\section{Шаардлагын мөшгөлтийн матриц (Requirement Traceability Matrix)}
Үндсэн функционал шаардлагууд туршилтаар хэрхэн баталгаажсаныг Хүснэгт \ref{table:rtm}-д нэгтгэн харуулав. Энэхүү матриц нь шаардлага бүрийг харгалзах туршилтын тохиолдолтой (Test Case) холбож, хэрэгжилтийг баталгаажуулдаг.

\begin{table}[H]
    \centering
    \caption{Шаардлагын мөшгөлтийн матриц (RTM)}
    \label{table:rtm}
    \resizebox{\textwidth}{!}{
        \begin{tabular}{|c|p{7cm}|c|c|}
            \hline
            \textbf{ID} & \textbf{Функционал шаардлага} & \textbf{Тест ID} & \textbf{Төлөв} \\
            \hline
            FR1 & Тогтмол илэрхийлэлд (Regex) суурилсан танилт & TC-001 & \checkmark Pass \\
            \hline
            FR2 & Хадгалсан харагдац (Хавтас) удирдлага & TC-002 & \checkmark Pass \\
            \hline
            FR3 & Нугалдаг төхөөрөмжийн (Foldable) дэмжлэг & TC-003 & \checkmark Pass \\
            \hline
            FR4 & Шудрах (Swipe) үйлдлийн тохиргоо & TC-004 & \checkmark Pass \\
            \hline
            FR5 & Хугацааны бүлэглэлт (Sticky Headers) & TC-005 & \checkmark Pass \\
            \hline
        \end{tabular}
    }
\end{table}
\section{Нэвтрүүлэлт ба Байршуулалт (Deployment \& Production)}

\subsection{Байршуулалтын орчин}
Энэхүү төсөл нь Android мобайл аппликейшн тул
байршуулалтын үйл явц нь APK үүсгэх, гарын үсэг зурах
(code signing) болон бодит төхөөрөмжид суулгах
үе шатуудаас бүрдэнэ.

\subsection{CI/CD Pipeline}
Хөгжүүлэлтийн эхний шатанд GitHub Actions ашиглан
автоматжуулсан CI/CD pipeline тохируулахыг оролдсон боловч
Fossify Messages-ийн суурь кодын нарийн төвөгтэй
build dependency-уудаас шалтгаалан pipeline тогтворгүй
ажиллах болсон тул энэ төслийн хүрээнд идэвхгүй болгосон.
Цаашдын хөгжүүлэлтэд CI/CD-г бүрэн нэвтрүүлэх нь
тэргүүлэх зорилтуудын нэг болно.

\subsection{Бодит төхөөрөмж дээрх туршилт}
Эцсийн APK хувилбарыг дараах техникийн үзүүлэлтүүдийг
хангасан байдлаар бүтээж, бодит төхөөрөмжүүд дээр
sideload аргаар суулгаж туршилт хийв.

\begin{itemize}
    \item \textbf{Target SDK}: Android 14 (API 34)
    \item \textbf{Minimum SDK}: Android 6.0 (API 23)
    \item \textbf{Build type}: Release (R8 оновчлолтой)
\end{itemize}

\begin{table}[H]
    \centering
    \caption{Бодит төхөөрөмж дээрх туршилтын үр дүн}
    \label{table:device_test}
    \resizebox{\textwidth}{!}{
        \begin{tabular}{|l|c|c|l|}
            \hline
            \textbf{Төхөөрөмж} &
            \textbf{Төрөл} &
            \textbf{Android} &
            \textbf{Үр дүн} \\
            \hline
            Samsung Galaxy Note20 Ultra &
            Ердийн дэлгэц &
            Android 13 &
            Pass \\
            \hline
            Samsung Galaxy Z Fold3 &
            Нугалдаг дэлгэц &
            Android 15 &
            Pass \\
            \hline
        \end{tabular}
    }
\end{table}

Нугалдаг дэлгэц бүхий төхөөрөмж болох Samsung Galaxy Z Fold3 дээр
аппликейшн нээгдэх болон хаагдах үед хоёр баганат
харагдац гацалтгүйгээр шилжиж байсныг бодит
орчинд баталгаажуулав. Ердийн дэлгэцтэй
Samsung Galaxy Note20 Ultra дээр бүх үндсэн
функцүүд зөв ажилласан.
\section{Эцсийн үнэлгээ}
Туршилтын үр дүнгээс үзэхэд систем нь хэрэглэгчийн тодорхойлсон дүрмийн дагуу зурвасыг алдаагүй ангилж, их хэмжээний өгөгдөл болон нугалдаг төхөөрөмжийн динамик нөхцөлд найдвартай ажиллах чадамжтай болох нь батлагдлаа.
